This is a template to make a thesis and/or slides in \href{https://www.latex-project.org}{\LaTeX}, free software to compile high-quality scientific documents. The template contains the typical parts such as a title page (for VUB, BRUFACE and ULB), acknowledgements, abstract, table of contents, a nomenclature, the typical chapters, appendices and a bibliography. In the appendices you will find practical information about formatting and examples of useful packages.

This template is too large to compile with a free Overleaf plan. Consult section A.1.4 in appendix A for possible workarounds. % a hard link to A.1.4 is added on purpose, because \cref will not work when appendix A is not included, as in myThesis.tex

A \href{https://vub.cloud.panopto.eu/Panopto/Pages/Sessions/List.aspx?folderID=eb5cbed1-a7f8-4d08-bca3-ad1e00db8d7a}{set of videos} is available that demonstrate the use of this template.

This template comes with four main files:
\begin{itemize}
    \item \texttt{mainReport}: to compile a report and all the appendices of this template
    \item \texttt{mainSlides}: to compile slides. Only selected content is put on the slides though. To see all content, you must compile \texttt{mainReport} and/or inspect the source code.
    \item \texttt{PhD}: to compile a PhD without the appendices of this template
    \item \texttt{Thesis}: to compile a thesis without the appendices of this template
\end{itemize}

The template uses the class \href{https://ctan.org/pkg/memoir?lang=en}{memoir}. It is a class for typesetting mathematical works as books, reports and articles.

This template is written in English. If you want to write in Dutch, you must select Dutch as the spell checking language in Overleaf via the Menu-button in the upper left corner and replace \texttt{english} by \texttt{dutch} as option in the \cs{documentclass} command.